\cleardoublepage
\chapter{Bibliografia}

A continuació es detallen les fonts acadèmiques i estudis científics consultats per a la fonamentació teòrica d'aquest treball, seguint la normativa APA:

\begin{itemize}
    \item Du, X., et al. (2020). \textit{How viable is password cracking in digital forensic investigation? Analyzing the guessability of over 3.9 billion real-world accounts}. Recuperat de: \url{https://forensicsandsecurity.com/papers/PasswordCracking3BillionAccounts.pdf}
    
    \item Wang, X., \& Ding, Y. (2023). \textit{No single silver bullet: Measuring the accuracy of password strength meters}. USENIX Security Symposium. Recuperat de: \url{https://www.usenix.org/system/files/usenixsecurity23-wang-ding-silver-bullet.pdf}
    
    \item Al-Shehri, A., et al. (2025). \textit{Analysing the impact of password length and complexity on the effectiveness of brute force attacks}. International Journal of Network Security \& Its Applications. Recuperat de: \url{https://aircconline.com/ijnsa/V17N2/17225ijnsa03.pdf}
    
    \item Gupta, S., et al. (2025). \textit{Optimizing password cracking for digital investigations}. arXiv. Recuperat de: \url{https://arxiv.org/pdf/2504.03347}
    
    \item Pratama, A., et al. (2025). \textit{Password strength study using the zxcvbn algorithm and brute-force time estimation to strengthen cybersecurity}. EJournal Universitas Nusa Mandiri. Recuperat de: \url{https://ejournal.nusamandiri.ac.id/index.php/pilar/article/download/6119/1361}
\end{itemize}

\vspace{1cm}
\small
\color{gray}
Nota: Totes les fonts han estat consultades durant el procés de redacció del treball (març de 2026).
